\documentclass{../../kalkulus-ppt}

\author[Tetew]{Teosofi Hidayah Agung}
\date{\today}
\title[Kalkulus 2 - Bab 3]{Integrasi Numerik dan Integrasi Tak Wajar}
\institute[Matematika ITS]{Departemen Matematika\\ Institut Teknologi Sepuluh Nopember}
\titlegraphic{{\includegraphics[scale=0.3]{ITS.png}$\quad$\includegraphics[scale=0.02]{M.png}}}

\setbeamertemplate{itemize items}[circle]
\pgfplotsset{
    grafik integral/.style={
        width=6.15cm,height=4.45cm,axis lines=middle,
        xlabel={$x$},ylabel={$y$},samples=100,
        ticklabel style={font=\scriptsize},label style={font=\scriptsize},
        every axis plot/.append style={thick}
    }
}

\begin{document}

\begin{frame}
    \titlepage
\end{frame}

\begin{frame}{Daftar isi}
    \tableofcontents
\end{frame}

% KALKULUS 2 BAB 3.pdf, halaman PDF 1--15 (cetak 107--121).
\section{Integrasi Numerik}

\begin{frame}{\insertsection}
    \framesubtitle{Mengganti luas dengan jumlah luas persegi panjang}
    \begin{columns}[c,onlytextwidth]
        \begin{column}{.53\textwidth}
            \begin{definisi}[Hampiran jumlah Riemann]
                Bagi $[a,b]$ menjadi $n$ subselang sama lebar
                $h=(b-a)/n$. Tulis $x_i=a+ih$. Untuk setiap
                $x_i^*\in[x_{i-1},x_i]$, gunakan
                \[
                    \int_a^b f(x)\,dx\approx h\sum_{i=1}^{n}f(x_i^*).
                \]
            \end{definisi}
            Setiap titik sampel menentukan tinggi untuk satu subselang.
            Jika $f$ kontinu, jumlah itu menuju integral saat $n\to\infty$.
            Contohnya, $\int_0^1 e^{-x^2}\,dx$ dapat dihampiri dari nilai-nilai $f$.
        \end{column}
        \begin{column}{.43\textwidth}
            \centering
            \begin{tikzpicture}
                \begin{axis}[grafik integral,width=5.85cm,clip mode=individual,
                    xmin=0,xmax=3.35,ymin=0,ymax=3.65,xtick={1,2,3},
                    xticklabels={$x_1$,$x_2$,$x_3$},ytick=\empty]
                    \filldraw[fill=HIMAmuda!30,draw=teal!70!black]
                        (axis cs:0,0) rectangle (axis cs:1,.8125);
                    \filldraw[fill=HIMAmuda!30,draw=teal!70!black]
                        (axis cs:1,0) rectangle (axis cs:2,1.3125);
                    \filldraw[fill=HIMAmuda!30,draw=teal!70!black]
                        (axis cs:2,0) rectangle (axis cs:3,2.3125);
                    \addplot[HIMAtua,very thick,domain=0:3] {.75+x^2/4};
                    \addplot[only marks,mark=*,mark size=1.6pt,HIMAtua]
                        coordinates {(.5,.8125) (1.5,1.3125) (2.5,2.3125)};
                    \node[font=\scriptsize,below=3pt] at (axis cs:0,0) {$x_0$};
                \end{axis}
            \end{tikzpicture}\par
            \small Tinggi persegi panjang diambil di titik tengah.
        \end{column}
    \end{columns}
\end{frame}

\begin{frame}{\insertsection}
    \framesubtitle{Tiga pilihan titik sampel pada partisi yang sama}
    \begin{block}{Ujung kiri, ujung kanan, dan titik tengah}
        Dengan $y_i=f(x_i)$ dan $m_i=(x_{i-1}+x_i)/2$,
        \[
            L_n=h\sum_{i=0}^{n-1}y_i,\qquad
            R_n=h\sum_{i=1}^{n}y_i,\qquad
            M_n=h\sum_{i=1}^{n} f(m_i).
        \]
    \end{block}
    \begin{contoh}
        Untuk fungsi yang naik, titik mana memberi tinggi terkecil dan terbesar
        di setiap subselang?
    \end{contoh}
    \onslide<2->{$f(x_{i-1})$ adalah tinggi terkecil dan $f(x_i)$ tinggi terbesar.
    Karena itu $L_n\le\int_a^b f(x)\,dx\le R_n$.
    Pada fungsi yang turun, urutan $L_n$ dan $R_n$ terbalik.}
\end{frame}

\begin{frame}{\insertsection}
    \framesubtitle{Rata-rata dua ujung membentuk trapesium}
    \begin{columns}[c,onlytextwidth]
        \begin{column}{.53\textwidth}
            \begin{teorema}[Aturan trapesium]
                \[
                    T_n=\tfrac12(L_n+R_n)
                        =\frac h2\left[y_0+2\sum_{i=1}^{n-1}y_i+y_n\right].
                \]
            \end{teorema}
            Pada $[x_{i-1},x_i]$, garis penghubung kedua ujung memberikan
            luas bertanda $\tfrac h2(y_{i-1}+y_i)$. Titik interior dipakai
            oleh dua trapesium, sehingga koefisiennya $2$.
        \end{column}
        \begin{column}{.43\textwidth}
            \centering
            \begin{tikzpicture}
                \begin{axis}[grafik integral,width=5.85cm,clip mode=individual,
                    xmin=0,xmax=3.35,ymin=0,ymax=3.65,xtick={1,2,3},
                    xticklabels={$x_1$,$x_2$,$x_3$},ytick=\empty]
                    \filldraw[fill=HIMAmuda!30,draw=teal!80!black]
                        (axis cs:0,0) -- (axis cs:0,.75)
                        -- (axis cs:1,1) -- (axis cs:1,0) -- cycle;
                    \filldraw[fill=HIMAmuda!22,draw=teal!80!black]
                        (axis cs:1,0) -- (axis cs:1,1)
                        -- (axis cs:2,1.75) -- (axis cs:2,0) -- cycle;
                    \filldraw[fill=HIMAmuda!30,draw=teal!80!black]
                        (axis cs:2,0) -- (axis cs:2,1.75)
                        -- (axis cs:3,3) -- (axis cs:3,0) -- cycle;
                    \addplot[HIMAtua,very thick,domain=0:3] {.75+x^2/4};
                    \node[font=\scriptsize,below=3pt] at (axis cs:0,0) {$x_0$};
                \end{axis}
            \end{tikzpicture}\par
            \small Tiap ruas hijau menghubungkan dua nilai ujung.
        \end{column}
    \end{columns}
\end{frame}

\begin{frame}{\insertsection}
    \framesubtitle{Luas garis singgung setara dengan hampiran titik tengah}
    \begin{columns}[c,onlytextwidth]
        \begin{column}{.52\textwidth}
            Di subselang selebar $h$ dengan titik tengah $m$, garis singgungnya
            $\ell(x)=f(m)+f'(m)(x-m)$.
            \[
                \int_{m-h/2}^{m+h/2}\ell(x)\,dx=h f(m).
            \]
            Suku $f'(m)(x-m)$ mempunyai integral nol karena kedua sisi $m$
            sama lebar. Jadi garis singgung dan garis datar $y=f(m)$
            memberi hampiran yang sama.
        \end{column}
        \begin{column}{.44\textwidth}
            \centering
            \begin{tikzpicture}
                \begin{axis}[grafik integral,width=5.95cm,xmin=0,xmax=2.25,
                    ymin=0,ymax=5.65,xtick={0,1,2},
                    xticklabels={$x_{i-1}$,$m$,$x_i$},ytick={2},yticklabels={$f(m)$}]
                    \fill[HIMAmuda!35] (axis cs:0,0) -- (axis cs:0,2) -- (axis cs:1,2) -- cycle;
                    \fill[HIMAmuda!35] (axis cs:1,2) -- (axis cs:2,2) -- (axis cs:2,4) -- cycle;
                    \addplot[HIMAtua,very thick,domain=0:2] {1+x^2};
                    \addplot[teal!85!black,domain=0:2] {2*x};
                    \draw[gray,dashed,thick] (axis cs:0,2) -- (axis cs:2,2);
                    \draw[gray,dashed] (axis cs:1,0) -- (axis cs:1,2);
                    \fill[HIMAtua] (axis cs:1,2) circle[radius=1.6pt];
                    \node[font=\scriptsize,HIMAtua,anchor=west] at (axis cs:.21,3.65) {$f(x)$};
                    \node[font=\scriptsize,teal!80!black,anchor=east] at (axis cs:1.98,3.52) {$\ell(x)$};
                \end{axis}
            \end{tikzpicture}\par
            \small Kedua segitiga berarsir memiliki luas yang sama.
        \end{column}
    \end{columns}
\end{frame}

\begin{frame}{\insertsection}
    \framesubtitle{Kecekungan menentukan arah galat}
    \begin{columns}[c,onlytextwidth]
        \begin{column}{.54\textwidth}
            \begin{teorema}[Posisi hampiran terhadap integral]
                Untuk $f$ kontinu dan $I=\int_a^b f(x)\,dx$,
                \[
                    \begin{array}{c|c}
                        \text{Bentuk kurva di }[a,b] & \text{Urutan}\\ \hline
                        \text{Cekung ke atas} & M_n\le I\le T_n\\
                        \text{Cekung ke bawah} & T_n\le I\le M_n
                    \end{array}
                \]
                Pada kecekungan tegas ketaksamaannya ketat.
            \end{teorema}
            Untuk cekung ke atas, tali busur berada di atas kurva.
            Jika $f$ terdiferensialkan, garis singgung berada di bawah kurva.
            Pada cekung ke bawah posisi keduanya terbalik.
        \end{column}
        \begin{column}{.42\textwidth}
            \centering
            \begin{tikzpicture}
                \begin{axis}[grafik integral,width=5.7cm,xmin=0,xmax=2.25,
                    ymin=0,ymax=3.65,xtick={0,1,2},xticklabels={$x_{i-1}$,$m$,$x_i$},
                    ytick=\empty]
                    \addplot[magenta,densely dashed,domain=0:2] {x+1};
                    \addplot[HIMAtua,very thick,domain=0:2] {1+x^2/2};
                    \addplot[teal!80!black,densely dashed,domain=0:2] {x+.5};
                    \addplot[only marks,mark=*,mark size=1.6pt,HIMAtua]
                        coordinates {(0,1) (1,1.5) (2,3)};
                    \node[font=\scriptsize,magenta,anchor=west] at (axis cs:.22,2.12) {tali busur};
                    \node[font=\scriptsize,teal!80!black,anchor=west] at (axis cs:.53,.61) {singgung};
                \end{axis}
            \end{tikzpicture}
        \end{column}
    \end{columns}
\end{frame}

\begin{frame}{\insertsection}
    \framesubtitle{Membandingkan ukuran galat titik tengah dan trapesium}
    \begin{definisi}[Galat bertanda]
        Untuk hampiran $A_n$ dari nilai eksak $I$, galatnya $E_A=I-A_n$.
        Besar galat adalah $|E_A|$. Dengan demikian $E_M=I-M_n$
        dan $E_T=I-T_n$.
    \end{definisi}
    \begin{teorema}[Perbandingan galat]
        Jika $f$ kontinu dan arah kecekungannya tetap pada $[a,b]$,
        maka $|E_M|\le |E_T|$.
    \end{teorema}
    \onslide<2->{Titik lama dan titik tengah menjadi titik ujung partisi baru.
    Karena itu $T_{2n}=(T_n+M_n)/2$.
    Untuk cekung ke atas, $I\le T_{2n}$, sehingga $I-M_n\le T_n-I$.
    Untuk cekung ke bawah, $I\ge T_{2n}$ dan tanda galat terbalik.}
    Uji arah kecekungan sebelum memakai perbandingan ini.
\end{frame}

\begin{frame}{\insertsection}
    \framesubtitle{Menghitung hampiran $\ln 3$ dengan empat subselang}
    \begin{columns}[c,onlytextwidth]
        \begin{column}{.39\textwidth}
            \centering
            \small
            \renewcommand{\arraystretch}{1.25}
            \begin{tabular}{c|c|c}
                $i$ & $y_i=1/x_i$ & $f(m_i)$\\ \hline
                $0$ & $1$ & ---\\
                $1$ & $2/3$ & $4/5$\\
                $2$ & $1/2$ & $4/7$\\
                $3$ & $2/5$ & $4/9$\\
                $4$ & $1/3$ & $4/11$
            \end{tabular}
        \end{column}
        \begin{column}{.57\textwidth}
            \begin{contoh}
                Hampiri $I=\int_1^3 dx/x=\ln 3$ memakai $h=(3-1)/4=1/2$.
                Titik tengah pertama $m_1=5/4$.
            \end{contoh}
            \onslide<2->{%
                \[
                    \begin{aligned}
                        M_4&=\tfrac12(\tfrac45+\tfrac47+\tfrac49+\tfrac4{11})
                            \approx 1{,}089755,\\
                        T_4&=\tfrac14[1+2(\tfrac23+\tfrac12+\tfrac25)+\tfrac13]
                            \approx 1{,}116667.
                    \end{aligned}
                \]
                Karena $f''(x)=2/x^3>0$, nilai eksak $\ln 3\approx1{,}098612$
                memang berada di antara keduanya.
            }
        \end{column}
    \end{columns}
\end{frame}

\begin{frame}{\insertsection}
    \framesubtitle{Simpson memakai satu parabola tiap dua subselang}
    \begin{columns}[c,onlytextwidth]
        \begin{column}{.51\textwidth}
            Hampiri $f(x)$ pada $[m-h,m+h]$ dengan polinomial derajat
            paling tinggi dua yang melalui tiga titik sampel.
            Tulis $p(m+t)=A+Bt+Ct^2$.
            \[
                \int_{-h}^{h}(A+Bt+Ct^2)\,dt=2Ah+\frac{2Ch^3}{3}.
            \]
            Jika $y_0=p(m-h)$, $y_1=p(m)$, dan $y_2=p(m+h)$, maka
            $y_0+4y_1+y_2=6A+2Ch^2$. Jadi luas bertandanya
            $\tfrac h3(y_0+4y_1+y_2)$.
        \end{column}
        \begin{column}{.45\textwidth}
            \centering
            \begin{tikzpicture}
                \begin{axis}[grafik integral,width=6.05cm,clip mode=individual,
                    xmin=0,xmax=2.25,ymin=0,ymax=8.4,xtick={1,2},
                    xticklabels={$m$,$m+h$},ytick=\empty]
                    \addplot[HIMAtua,very thick,domain=0:2] {exp(x)};
                    \addplot[teal!80!black,dashed,domain=0:2] {1+.242036*x+1.476246*x^2};
                    \draw[gray,dotted] (axis cs:1,0) -- (axis cs:1,2.718282);
                    \draw[gray,dotted] (axis cs:2,0) -- (axis cs:2,7.389056);
                    \addplot[only marks,mark=*,mark size=1.7pt,HIMAtua]
                        coordinates {(0,1) (1,2.718282) (2,7.389056)};
                    \node[font=\scriptsize,below=3pt] at (axis cs:0,0) {$m-h$};
                    \node[font=\scriptsize,HIMAtua,anchor=west] at (axis cs:.27,6.8) {kurva $f$};
                    \node[font=\scriptsize,teal!80!black,anchor=west] at (axis cs:.27,5.55) {parabola $p$};
                \end{axis}
            \end{tikzpicture}\par
            \small Ketiga titik sampel terletak pada $f$ dan $p$.
        \end{column}
    \end{columns}
\end{frame}

\begin{frame}{\insertsection}
    \framesubtitle{Bobot Simpson berasal dari pasangan subselang}
    \begin{teorema}[Aturan Simpson]
        Untuk $n$ genap dan $h=(b-a)/n$,
        \[
            S_n=\frac h3\left[y_0+4y_1+2y_2+4y_3+\cdots+
                            2y_{n-2}+4y_{n-1}+y_n\right].
        \]
        Untuk $n=2$, gunakan satu pasangan dengan bobot $1,4,1$.
    \end{teorema}
    \onslide<2->{%
        Dua pasangan pertama menyumbang
        \[
            \underbrace{\tfrac h3(y_0+4y_1+y_2)}_{\text{pasangan pertama}}+
            \underbrace{\tfrac h3(y_2+4y_3+y_4)}_{\text{pasangan kedua}}.
        \]
        Titik interior genap seperti $x_2$ menjadi batas dua pasangan.
        Titik berindeks ganjil menjadi titik tengah satu pasangan.
    }
\end{frame}

\begin{frame}{\insertsection}
    \framesubtitle{Aturan Simpson pada data $\ln 3$ yang sama}
    \begin{contoh}
        Gunakan kembali $y_0,\ldots,y_4$ dari $I=\int_1^3 dx/x$.
        Kali ini pasangkan empat subselang untuk memakai Simpson.
    \end{contoh}
    \onslide<2->{%
        \[
            \begin{aligned}
                S_4
                &=\frac16\left[1+4(\tfrac23)+2(\tfrac12)+4(\tfrac25)+\tfrac13\right]\\
                &=\frac{11}{10}=1{,}1.
            \end{aligned}
        \]
    }
    \onslide<3->{%
        \begin{block}{Seberapa dekat hasilnya?}
            $|E_S|=|\ln3-1{,}1|\approx0{,}001388$.
            Pada contoh ini, galat Simpson lebih kecil daripada galat titik tengah
            dan trapesium walaupun menggunakan partisi yang sama.
        \end{block}
    }
\end{frame}

\begin{frame}{\insertsection}
    \framesubtitle{Dua sumber galat saat menghitung secara numerik}
    \begin{columns}[t,onlytextwidth]
        \begin{column}{.48\textwidth}
            \begin{block}{Galat pemotongan}
                $I-A_n$ muncul karena rumus memakai sejumlah hingga sampel.
                Galat ini tetap ada walaupun seluruh operasi dilakukan eksak.
            \end{block}
        \end{column}
        \begin{column}{.48\textwidth}
            \begin{block}{Galat pembulatan}
                $A_n-\widetilde A_n$ muncul ketika nilai sampel atau
                hasil operasi ditulis dengan digit terbatas.
            \end{block}
        \end{column}
    \end{columns}
    \[
        I-\widetilde A_n=(I-A_n)+(A_n-\widetilde A_n).
    \]
    \onslide<2->{Batas galat berikut mengukur bagian $I-A_n$.
    Perhitungan dengan banyak subselang dapat memerlukan lebih banyak digit
    agar pembulatan tetap kecil.}
\end{frame}

\begin{frame}{\insertsection}
    \framesubtitle{Batas galat dari turunan fungsi}
    \begin{columns}[t,onlytextwidth]
        \begin{column}{.50\textwidth}
            \begin{teorema}[Titik tengah dan trapesium]
                Jika $f''$ kontinu dan $|f''(x)|\le K_2$ pada $[a,b]$, maka
                \[
                    \begin{aligned}
                        |E_M|&\le\frac{(b-a)^3K_2}{24n^2},\\
                        |E_T|&\le\frac{(b-a)^3K_2}{12n^2}.
                    \end{aligned}
                \]
            \end{teorema}
        \end{column}
        \begin{column}{.46\textwidth}
            \begin{teorema}[Simpson]
                Jika $f^{(4)}$ kontinu dan $|f^{(4)}(x)|\le K_4$ pada $[a,b]$,
                maka untuk $n$ genap,
                \[
                    |E_S|\le\frac{(b-a)^5K_4}{180n^4}.
                \]
            \end{teorema}
        \end{column}
    \end{columns}
    Gunakan konstanta $K_2$ atau $K_4$ yang membatasi turunan di seluruh selang.
    Batas itu menjamin toleransi sekalipun nilai eksak $I$ belum diketahui.
\end{frame}

\begin{frame}{\insertsection}
    \framesubtitle{Membandingkan batas dan galat aktual pada contoh yang sama}
    Untuk $f(x)=1/x$ di $[1,3]$, berlaku $f''(x)=2/x^3$ dan
    $f^{(4)}(x)=24/x^5$. Keduanya maksimum di $x=1$, jadi ambil
    $K_2=2$ dan $K_4=24$.
    \onslide<2->{%
        \begin{center}
            \renewcommand{\arraystretch}{1.25}
            \begin{tabular}{l|c|c|c}
                \textbf{Metode} & \textbf{Hampiran} & \textbf{Besar galat aktual}
                                & \textbf{Batas atas}\\ \hline
                Titik tengah & $1{,}089755$ & $\approx 0{,}008858$ & $1/24$\\
                Trapesium    & $1{,}116667$ & $\approx 0{,}018054$ & $1/12$\\
                Simpson      & $1{,}100000$ & $\approx 0{,}001388$ & $1/60$
            \end{tabular}
        \end{center}
        Semua galat aktual, dihitung terhadap $\ln 3$, berada di bawah batasnya.
        Misalnya batas Simpson adalah
        $(3-1)^5(24)/(180\cdot4^4)=1/60$.
    }
\end{frame}

\begin{frame}{\insertsection}
    \framesubtitle{Memilih jumlah subselang dari toleransi}
    \begin{contoh}
        Berapa $n$ genap terkecil yang dijamin oleh batas galat Simpson
        agar hampiran $\int_1^3 dx/x$ memiliki besar galat $<10^{-6}$?
    \end{contoh}
    \onslide<2->{Untuk $K_4=24$ dan $b-a=2$, syarat yang cukup ialah
    \[
        \frac{2^5(24)}{180n^4}<10^{-6}
        \quad\Longleftrightarrow\quad
        n>\left(\frac{64\cdot10^6}{15}\right)^{1/4}
            \approx45{,}449.
    \]}
    \onslide<3->{Bilangan genap terkecil yang melewati ambang ini adalah
    $n=46$. Pilihan ini cukup untuk menjamin toleransi dari batas tersebut.}
\end{frame}

\begin{frame}{\insertsection}
    \framesubtitle{Batas yang mudah dihitung sudah dapat dipakai}
    \begin{contoh}
        Pilih $n$ agar aturan titik tengah untuk $\int_0^1\cos(x^2)\,dx$
        memiliki besar galat $<2\cdot10^{-4}$.
    \end{contoh}
    \onslide<2->{%
        \[
            f''(x)=-2\sin(x^2)-4x^2\cos(x^2),\qquad
            |f''(x)|\le 2+4x^2\le6.
        \]
        Maka $K_2=6$ sah sebagai batas atas, dengan
        $|E_M|\le 6/(24n^2)=1/(4n^2)$.
    }
    \onslide<3->{Syarat $1/(4n^2)<2\cdot10^{-4}$ memberi
    $n>\sqrt{1250}\approx35{,}355$, jadi $n=36$ cukup.
    Ketaksamaan segitiga menghindari pencarian maksimum eksak $|f''|$.}
\end{frame}

\begin{frame}{\insertsection}
    \framesubtitle{Pengaruh membagi dua lebar subselang}
    \begin{columns}[c,onlytextwidth]
        \begin{column}{.48\textwidth}
            \begin{block}{Tulis kembali dengan $h=(b-a)/n$}
                \[
                    \begin{aligned}
                        |E_M|&\le\tfrac{K_2(b-a)}{24}h^2,\\
                        |E_T|&\le\tfrac{K_2(b-a)}{12}h^2,\\
                        |E_S|&\le\tfrac{K_4(b-a)}{180}h^4.
                    \end{aligned}
                \]
            \end{block}
        \end{column}
        \begin{column}{.48\textwidth}
            Jika $h$ menjadi $h/2$, batas titik tengah dan trapesium
            menjadi $1/4$ kali semula. Batas Simpson menjadi $1/16$ kali semula.
            \onslide<2->{%
                \begin{block}{Kapan hasilnya eksak?}
                    Untuk polinomial berderajat paling tinggi satu,
                    $K_2=0$ membuat $M_n$ dan $T_n$ eksak.
                    Untuk derajat paling tinggi tiga, $K_4=0$ membuat $S_n$ eksak.
                \end{block}
            }
        \end{column}
    \end{columns}
\end{frame}

\begin{frame}{\insertsection}
    \framesubtitle{Menghampiri jarak dari data kecepatan}
    \begin{columns}[c,onlytextwidth]
        \begin{column}{.53\textwidth}
            \begin{contoh}
                Kecepatan benda diukur tiap $2$ detik. Pada $t=0,2,4,6,8$
                detik, hasil ukurnya $0,3,5,6,6{,}5$ m/s.
                Hampiri jarak selama $8$ detik.
            \end{contoh}
            \onslide<2->{Karena kecepatannya nonnegatif, jarak dihampiri dengan
            Simpson pada $\int_0^8 v(t)\,dt$.
            \[
                \tfrac23[0+4(3)+2(5)+4(6)+6{,}5]=35\text{ m}.
            \]
            Lima titik ukur membentuk empat subselang.}
        \end{column}
        \begin{column}{.43\textwidth}
            \centering
            \begin{tikzpicture}
                \begin{axis}[grafik integral,width=5.8cm,xlabel={$t$ (s)},ylabel={$v$ (m/s)},
                    xmin=0,xmax=9.2,ymin=0,ymax=7.8,xtick={0,2,4,6,8},ytick={3,5,6.5}]
                    \draw[gray,dotted] (axis cs:2,0) -- (axis cs:2,3);
                    \draw[gray,dotted] (axis cs:4,0) -- (axis cs:4,5);
                    \draw[gray,dotted] (axis cs:6,0) -- (axis cs:6,6);
                    \draw[gray,dotted] (axis cs:8,0) -- (axis cs:8,6.5);
                    \addplot[only marks,mark=*,mark size=2pt,HIMAtua]
                        coordinates {(0,0) (2,3) (4,5) (6,6) (8,6.5)};
                \end{axis}
            \end{tikzpicture}
        \end{column}
    \end{columns}
\end{frame}

\begin{frame}{\insertsection}
    \framesubtitle{Latihan integrasi numerik}
    \begin{latihan}
        Gunakan $I=\int_0^2 x^2\,dx$.
        \begin{enumerate}
            \item Hitung $M_4$ dan $T_4$. Di mana posisi $I$?
            \item Tentukan $S_4$ dan jelaskan mengapa hasilnya eksak.
            \item Pilih $n$ agar batas galat titik tengah kurang dari $10^{-3}$.
        \end{enumerate}
    \end{latihan}
    \onslide<2->{%
        \begin{block}{Cek jawaban}
            \begin{enumerate}
                \item $M_4=21/8 < I=8/3 < T_4=11/4$ karena $f''=2>0$.
                \item $S_4=8/3$ karena turunan keempat $x^2$ nol.
                \item $|E_M|\le 2/(3n^2)$. Ambil $n=26$ agar batasnya $<10^{-3}$.
            \end{enumerate}
        \end{block}
    }
\end{frame}

% KALKULUS 2 BAB 3.pdf, halaman PDF 15--24 (cetak 121--130).
\section{Integral Tak Wajar}

\begin{frame}{\insertsection}
    \framesubtitle{Kapan integral perlu dibaca sebagai limit?}
    \begin{definisi}[Integral tak wajar]
        Integral tak wajar didefinisikan melalui limit integral pada subselang
        berhingga yang menghindari titik singular. Integral itu \textbf{konvergen}
        jika setiap limit yang diperlukan ada dan berhingga.
        Jika ada limit yang tak berhingga atau tidak ada, integralnya \textbf{divergen}.
    \end{definisi}
    \begin{columns}[t,onlytextwidth]
        \begin{column}{.48\textwidth}
            \begin{contoh}
                $\int_2^\infty x^{-2}\,dx$ mempunyai batas atas tak berhingga.
                Potong ekornya di $R$ lalu besarkan $R$.
            \end{contoh}
        \end{column}
        \begin{column}{.48\textwidth}
            \begin{contoh}
                $\int_0^1 x^{-1/2}\,dx$ mempunyai integran tak terbatas di $0$.
                Mulai dari $s>0$ lalu dekatkan $s$ ke $0$.
            \end{contoh}
        \end{column}
    \end{columns}
\end{frame}

\begin{frame}{\insertsection}
    \framesubtitle{Batas atas bergerak menuju tak hingga}
    \begin{columns}[c,onlytextwidth]
        \begin{column}{.52\textwidth}
            \begin{definisi}[{Selang $[a,\infty)$}]
                Jika $f$ kontinu pada $[a,\infty)$, tetapkan
                \[
                    \int_a^\infty f(x)\,dx
                    =\lim_{R\to\infty}\int_a^R f(x)\,dx
                \]
                ketika limitnya berhingga.
            \end{definisi}
            Untuk $f\ge0$, bagian berarsir memberi luas hingga $R$.
            Ketika $R$ dinaikkan, luasnya bertambah dan mungkin mendekati
            suatu nilai hingga.
        \end{column}
        \begin{column}{.44\textwidth}
            \centering
            \begin{tikzpicture}
                \begin{axis}[grafik integral,width=5.9cm,xmin=0,xmax=7.6,
                    ymin=0,ymax=.34,xtick={2,4.45},xticklabels={$a$,$R$},ytick=\empty]
                    \addplot[draw=none,fill=HIMAmuda!36,domain=2:4.45] {1/x^2} \closedcycle;
                    \addplot[HIMAtua,very thick,domain=2:7.4] {1/x^2};
                    \draw[gray,dashed] (axis cs:4.45,0) -- (axis cs:4.45,.0505);
                    \draw[->,teal!80!black] (axis cs:4.45,.1) -- (axis cs:6.75,.1);
                    \node[font=\scriptsize,teal!80!black,anchor=south]
                        at (axis cs:5.6,.1) {$R\to\infty$};
                    \node[font=\scriptsize,HIMAtua,anchor=west] at (axis cs:2.55,.24) {$y=1/x^2$};
                \end{axis}
            \end{tikzpicture}
        \end{column}
    \end{columns}
\end{frame}

\begin{frame}{\insertsection}
    \framesubtitle{Dua kurva menuju nol dengan luas berbeda}
    \begin{columns}[t,onlytextwidth]
        \begin{column}{.48\textwidth}
            \begin{contoh}
                $\displaystyle\int_2^\infty \frac{dx}{x^2}$
            \end{contoh}
            \onslide<2->{%
                \[
                    \lim_{R\to\infty}\left[-\frac1x\right]_2^R
                    =\lim_{R\to\infty}\left(\frac12-\frac1R\right)
                    =\frac12.
                \]
                Luas totalnya berhingga.}
        \end{column}
        \begin{column}{.48\textwidth}
            \begin{contoh}
                $\displaystyle\int_2^\infty \frac{dx}{x}$
            \end{contoh}
            \onslide<3->{%
                \[
                    \lim_{R\to\infty}[\ln x]_2^R
                    =\lim_{R\to\infty}(\ln R-\ln2)=+\infty.
                \]
                Integral divergen karena luas terus bertambah tanpa batas.}
        \end{column}
    \end{columns}
    \onslide<3->{Kedua integran menuju nol. Laju penurunan menentukan
    apakah luas yang terakumulasi berhingga.}
\end{frame}

\begin{frame}{\insertsection}
    \framesubtitle{Batas bawah bergerak menuju negatif tak hingga}
    \begin{columns}[c,onlytextwidth]
        \begin{column}{.53\textwidth}
            \begin{definisi}[{Selang $(-\infty,b]$}]
                Jika $f$ kontinu pada $(-\infty,b]$,
                \[
                    \int_{-\infty}^{b}f(x)\,dx
                       =\lim_{L\to-\infty}\int_L^{b}f(x)\,dx
                \]
                ketika limitnya berhingga.
            \end{definisi}
            \begin{contoh}
                Karena $e^{2L}\to0$ ketika $L\to-\infty$,
                $\int_{-\infty}^1e^{2x}\,dx=
                \lim_{L\to-\infty}\tfrac12(e^2-e^{2L})=e^2/2$.
            \end{contoh}
        \end{column}
        \begin{column}{.43\textwidth}
            \centering
            \begin{tikzpicture}
                \begin{axis}[grafik integral,width=5.85cm,xmin=-3.25,xmax=1.35,
                    ymin=0,ymax=8.6,xtick={-1.3,1},xticklabels={$L$,$b$},ytick=\empty]
                    \addplot[draw=none,fill=HIMAmuda!36,domain=-1.3:1] {exp(2*x)} \closedcycle;
                    \addplot[HIMAtua,very thick,domain=-3.05:1] {exp(2*x)};
                    \draw[<-,teal!80!black] (axis cs:-2.9,.9) -- (axis cs:-1.3,.9);
                    \node[font=\scriptsize,teal!80!black,anchor=south]
                        at (axis cs:-2.1,.9) {$L\to-\infty$};
                    \draw[gray,dotted] (axis cs:1,0) -- (axis cs:1,7.389056);
                \end{axis}
            \end{tikzpicture}
        \end{column}
    \end{columns}
\end{frame}

\begin{frame}{\insertsection}
    \framesubtitle{Dua ujung tak hingga memerlukan dua limit}
    \begin{definisi}[Selang $(-\infty,\infty)$]
        Untuk suatu titik pemisah $c\in\R$,
        \[
            \int_{-\infty}^\infty f(x)\,dx
              =\int_{-\infty}^{c}f(x)\,dx+\int_c^\infty f(x)\,dx
        \]
        jika kedua integral di kanan konvergen. Titik $c$ dapat dipilih
        agar perhitungan mudah, tanpa mengubah hasil.
    \end{definisi}
    \begin{contoh}
        Untuk $f(x)=1/(4+x^2)$, pilih $c=0$.
        \onslide<2->{Anti-turunannya $\tfrac12\arctan(x/2)$.
        Integral di setiap sisi adalah $\pi/4$, sehingga totalnya $\pi/2$.}
    \end{contoh}
\end{frame}

\begin{frame}{\insertsection}
    \framesubtitle{Apa yang terjadi pada pemotongan simetris?}
    \begin{columns}[c,onlytextwidth]
        \begin{column}{.52\textwidth}
            \begin{contoh}
                Untuk $f(x)=x$, dua potongan simetris memberi
                $\int_{-R}^R x\,dx=0$ untuk setiap $R>0$.
            \end{contoh}
            \onslide<2->{Dengan titik pemisah $0$, dua limit yang diperlukan ialah
            \[
                \begin{aligned}
                    \lim_{L\to-\infty}\int_L^0 x\,dx&=-\infty,\\
                    \lim_{R\to\infty}\int_0^R x\,dx&=+\infty.
                \end{aligned}
            \]
            Jadi integral pada seluruh garis real divergen.}
        \end{column}
        \begin{column}{.44\textwidth}
            \centering
            \begin{tikzpicture}
                \begin{axis}[grafik integral,width=5.85cm,axis equal image,
                    xmin=-2.25,xmax=2.25,ymin=-2.25,ymax=2.25,
                    xtick={-1.7,1.7},xticklabels={$-R$,$R$},ytick=\empty]
                    \filldraw[fill=magenta!20,draw=magenta!80!black]
                        (axis cs:-1.7,0) -- (axis cs:-1.7,-1.7) -- (axis cs:0,0) -- cycle;
                    \filldraw[fill=HIMAmuda!35,draw=teal!80!black]
                        (axis cs:0,0) -- (axis cs:1.7,0) -- (axis cs:1.7,1.7) -- cycle;
                    \addplot[HIMAtua,very thick,domain=-2:2] {x};
                    \node[font=\scriptsize,magenta!80!black] at (axis cs:-1.29,-.45) {$-$};
                    \node[font=\scriptsize,teal!80!black] at (axis cs:1.29,.45) {$+$};
                \end{axis}
            \end{tikzpicture}\par
            \small Luas bertanda saling hapus pada setiap potongan simetris.
        \end{column}
    \end{columns}
\end{frame}

\begin{frame}{\insertsection}
    \framesubtitle{Dekati titik singular dari dalam selang}
    \begin{columns}[c,onlytextwidth]
        \begin{column}{.54\textwidth}
            \begin{definisi}[Titik singular di ujung]
                Jika $f$ kontinu kecuali di ujung yang ditunjukkan, gunakan
                \[
                    \begin{aligned}
                        \text{di }a\colon\quad
                        \int_a^b f(x)\,dx&=\lim_{s\to a^+}\int_s^b f(x)\,dx,\\
                        \text{di }b\colon\quad
                        \int_a^b f(x)\,dx&=\lim_{t\to b^-}\int_a^t f(x)\,dx.
                    \end{aligned}
                \]
                Integral konvergen jika limit yang sesuai berhingga.
            \end{definisi}
            Pada gambar, $t$ mendekati $b$ dari kiri. Integral hingga $t$
            masih berupa integral pada selang biasa.
        \end{column}
        \begin{column}{.42\textwidth}
            \centering
            \begin{tikzpicture}
                \begin{axis}[grafik integral,width=5.75cm,clip mode=individual,
                    xmin=0,xmax=4.58,ymin=0,ymax=2.61,xtick={2.9,4},xticklabels={$t$,$b$},
                    ytick=\empty]
                    \addplot[draw=none,fill=HIMAmuda!36,domain=0:2.9] {1/sqrt(4-x)} \closedcycle;
                    \addplot[HIMAtua,very thick,domain=0:3.86] {1/sqrt(4-x)};
                    \draw[gray,dashed] (axis cs:4,0) -- (axis cs:4,2.48);
                    \draw[gray,dotted] (axis cs:2.9,0) -- (axis cs:2.9,.95346);
                    \node[font=\scriptsize,below=3pt] at (axis cs:0,0) {$a$};
                    \node[font=\scriptsize,HIMAtua,anchor=west] at (axis cs:.31,1.75)
                        {$y=(4-x)^{-1/2}$};
                \end{axis}
            \end{tikzpicture}
        \end{column}
    \end{columns}
\end{frame}

\begin{frame}{\insertsection}
    \framesubtitle{Pangkat akar dapat memberi luas yang berhingga}
    \begin{columns}[t,onlytextwidth]
        \begin{column}{.48\textwidth}
            \begin{contoh}
                $\displaystyle\int_0^4 \frac{dx}{\sqrt{4-x}}$
            \end{contoh}
            \onslide<2->{Potong sebelum $x=4$.
            \[
                \lim_{t\to4^-}\bigl[4-2\sqrt{4-t}\bigr]=4.
            \]
            Integral konvergen walaupun kurvanya meninggi tanpa batas.}
        \end{column}
        \begin{column}{.48\textwidth}
            \begin{contoh}
                $\displaystyle\int_0^1-\frac1x\,dx$
            \end{contoh}
            \onslide<3->{Mulai dari $s>0$.
            \[
                \lim_{s\to0^+}[-\ln x]_s^1
                    =\lim_{s\to0^+}\ln s=-\infty.
            \]
            Integral divergen karena limitnya tak berhingga.}
        \end{column}
    \end{columns}
\end{frame}

\begin{frame}{\insertsection}
    \framesubtitle{Titik singular di bagian dalam menjadi batas baru}
    \begin{columns}[c,onlytextwidth]
        \begin{column}{.54\textwidth}
            \begin{definisi}[Titik singular $a<c<b$]
                Jika $f$ kontinu pada $[a,b]$ kecuali di $c$, tetapkan
                \[
                    \int_a^b f(x)\,dx
                    =\lim_{t\to c^-}\int_a^t f(x)\,dx+
                     \lim_{s\to c^+}\int_s^b f(x)\,dx
                \]
                apabila kedua limit berhingga.
            \end{definisi}
            Dari kiri, batas atas mendekati $c$. Dari kanan, batas bawah
            mendekati $c$. Jika ada beberapa titik singular, pisahkan
            selang sehingga setiap limit satu sisi dapat diperiksa.
        \end{column}
        \begin{column}{.42\textwidth}
            \centering
            \begin{tikzpicture}
                \begin{axis}[grafik integral,width=5.75cm,xmin=0,xmax=4.43,
                    ymin=0,ymax=4.44,xtick={0,1,4},xticklabels={$a$,$c$,$b$},
                    ytick=\empty]
                    \addplot[draw=none,fill=HIMAmuda!35,domain=0:.6]
                        {1/(1-x)^(2/3)} \closedcycle;
                    \addplot[draw=none,fill=PastelGreen!38,domain=1.42:4]
                        {1/(x-1)^(2/3)} \closedcycle;
                    \addplot[HIMAtua,very thick,domain=0:.892] {1/(1-x)^(2/3)};
                    \addplot[HIMAtua,very thick,domain=1.108:4] {1/(x-1)^(2/3)};
                    \draw[gray,dashed] (axis cs:1,0) -- (axis cs:1,4.34);
                    \draw[->,teal!80!black] (axis cs:.6,2.65) -- (axis cs:.93,2.65);
                    \draw[->,teal!80!black] (axis cs:1.42,2.65) -- (axis cs:1.07,2.65);
                \end{axis}
            \end{tikzpicture}
        \end{column}
    \end{columns}
\end{frame}

\begin{frame}{\insertsection}
    \framesubtitle{Hitung kedua sisi secara terpisah}
    \begin{contoh}
        Tentukan $\int_0^9 |x-1|^{-2/3}\,dx$.
    \end{contoh}
    \onslide<2->{Pisahkan di $x=1$. Di sebelah kiri ambil $u=1-x$,
    dan di sebelah kanan ambil $u=x-1$.
    \[
        \begin{aligned}
            \int_0^1 \frac{dx}{(1-x)^{2/3}}
              &=\lim_{\varepsilon\to0^+}[3u^{1/3}]_\varepsilon^1=3,\\
            \int_1^9 \frac{dx}{(x-1)^{2/3}}
              &=\lim_{\varepsilon\to0^+}[3u^{1/3}]_\varepsilon^8=6.
        \end{aligned}
    \]}
    \onslide<3->{Kedua limit konvergen. Jadi nilai integralnya $3+6=9$.}
\end{frame}

\begin{frame}{\insertsection}
    \framesubtitle{Periksa titik singular sebelum memakai anti-turunan}
    \begin{contoh}
        Evaluasi $\int_{-2}^2 x^{-2}\,dx$. Apa yang terjadi jika
        $\left[-1/x\right]_{-2}^2$ langsung dihitung?
    \end{contoh}
    \onslide<2->{Substitusi dua batas langsung memberi $-1$,
    padahal integrannya positif. Teorema Fundamental Kalkulus memerlukan
    perhitungan pada selang yang tidak melintasi $0$.}
    \onslide<3->{Potong di sebelah kanan $0$, misalnya.
    \[
        \int_0^2 \frac{dx}{x^2}
        =\lim_{s\to0^+}\left(\frac1s-\frac12\right)=+\infty.
    \]
    Satu sisi saja sudah divergen, sehingga integral pada $[-2,2]$ divergen.}
\end{frame}

\begin{frame}{\insertsection}
    \framesubtitle{Uji pangkat pada tak hingga dan pada nol}
    \begin{teorema}[Integral $p$]
        Untuk $p\in\R$,
        \[
            \begin{array}{c|c|c}
                \text{Integral} & \text{Konvergen jika} & \text{Nilai saat konvergen}\\ \hline
                \displaystyle\int_1^\infty x^{-p}\,dx & p>1 & 1/(p-1)\\[4pt]
                \displaystyle\int_0^1 x^{-p}\,dx & p<1 & 1/(1-p)
            \end{array}
        \]
    \end{teorema}
    \onslide<2->{Untuk $p\ne1$, anti-turunannya $x^{1-p}/(1-p)$.
    Suku $R^{1-p}$ menuju nol jika $p>1$, sedangkan $s^{1-p}$ menuju nol
    saat $s\to0^+$ jika $p<1$. Pada $p=1$ kedua integral divergen
    karena anti-turunannya $\ln x$.}
\end{frame}

\begin{frame}{\insertsection}
    \framesubtitle{Gunakan perbandingan untuk integran nonnegatif}
    \begin{teorema}[Uji perbandingan langsung]
        Jika $0\le f(x)\le g(x)$ untuk $x\ge a$ dan kedua fungsi kontinu, maka
        \begin{itemize}
            \item Konvergensi $\int_a^\infty g$ menjamin konvergensi $\int_a^\infty f$.
            \item Divergensi $\int_a^\infty f$ menjamin divergensi $\int_a^\infty g$.
        \end{itemize}
    \end{teorema}
    \onslide<2->{Integral potongan dari $f$ dan $g$ naik saat batas atas dinaikkan.
    Jika integral potongan yang lebih besar tetap terbatas, keduanya punya
    limit berhingga.}
    \onslide<3->{%
        \begin{contoh}
            Karena $0<1/(x^2+5)\le 1/x^2$ untuk $x\ge1$,
            $\int_1^\infty dx/(x^2+5)$ konvergen dan nilainya paling besar $1$.
        \end{contoh}
    }
\end{frame}

\begin{frame}{\insertsection}
    \framesubtitle{Perbandingan dari bawah dapat membuktikan divergensi}
    \begin{contoh}
        Tentukan konvergensi $\displaystyle\int_1^\infty\frac{dx}{\sqrt{x^2+3}}$.
    \end{contoh}
    \onslide<2->{Untuk $x\ge1$, berlaku $x^2+3\le4x^2$.
    Setelah mengambil akar dan kebalikan,
    \[
        \frac{1}{\sqrt{x^2+3}}\ge\frac1{2x}.
    \]}
    \onslide<3->{Integral $\tfrac12\int_1^\infty dx/x$ divergen.
    Maka integral yang ditanyakan juga divergen.
    Pembanding yang divergen dipakai sebagai batas bawah.}
\end{frame}

\begin{frame}{\insertsection}
    \framesubtitle{Potong ekor sebelum menghitung secara numerik}
    \begin{columns}[c,onlytextwidth]
        \begin{column}{.54\textwidth}
            \begin{contoh}
                Hampiri $\int_0^\infty e^{-x^2}\,dx$ dengan memotong di $x=3$
                dan menerapkan Simpson pada $[0,3]$.
            \end{contoh}
            \onslide<2->{Dengan $n=12$ dan $h=0{,}25$, diperoleh
            $S_{12}\approx0{,}886207$ pada $[0,3]$.}
            \onslide<3->{Untuk $x\ge3$, $e^{-x^2}\le (x/3)e^{-x^2}$.
            Maka ekor yang dipotong memenuhi
            \[
                0<\int_3^\infty e^{-x^2}\,dx
                  \le\frac{e^{-9}}6\approx2{,}057\cdot10^{-5}.
            \]
            Galat Simpson dan pembulatan pada $[0,3]$ dihitung tersendiri.}
        \end{column}
        \begin{column}{.42\textwidth}
            \centering
            \begin{tikzpicture}
                \begin{axis}[grafik integral,width=5.75cm,xmin=0,xmax=3.8,
                    ymin=0,ymax=1.16,xtick={0,1,2,3},ytick={1}]
                    \addplot[draw=none,fill=HIMAmuda!35,domain=0:3]
                        {exp(-(x^2))} \closedcycle;
                    \addplot[HIMAtua,very thick,domain=0:3.7] {exp(-(x^2))};
                    \draw[magenta,dashed] (axis cs:3,0) -- (axis cs:3,.6);
                    \node[font=\scriptsize,magenta,anchor=east] at (axis cs:2.92,.48) {batas potong};
                    \node[font=\scriptsize,HIMAtua,anchor=west] at (axis cs:.72,.77)
                        {$y=e^{-x^2}$};
                \end{axis}
            \end{tikzpicture}\par
            \small Integral eksaknya $\sqrt{\pi}/2\approx0{,}886227$.
        \end{column}
    \end{columns}
\end{frame}

\begin{frame}{\insertsection}
    \framesubtitle{Syarat konvergensi pada transformasi Laplace}
    \begin{definisi}[Transformasi Laplace]
        Untuk $s\in\R$ yang membuat integralnya konvergen,
        \[
            \mathcal L\{f\}(s)=\int_0^\infty e^{-st}f(t)\,dt.
        \]
    \end{definisi}
    \begin{contoh}
        Tentukan transformasi $f(t)=e^{3t}$ beserta syarat pada $s$.
    \end{contoh}
    \onslide<2->{Integrannya $e^{-(s-3)t}$. Jika $s>3$, eksponensial meluruh dan
    $\mathcal L\{e^{3t}\}(s)=1/(s-3)$.
    Jika $s=3$, integrannya $1$. Jika $s<3$, integrannya tumbuh.
    Kedua kasus terakhir memberi integral divergen.}
\end{frame}

\begin{frame}{\insertsection}
    \framesubtitle{Latihan integral tak wajar}
    \begin{latihan}
        Tentukan konvergensi. Hitung nilai integral jika konvergen.
        \begin{enumerate}
            \item $\int_0^\infty e^{-3x}\,dx$.
            \item $\int_0^2 (x-1)^{-1}\,dx$.
            \item $\int_0^\infty x^{-p}\,dx$ untuk $p\in\R$.
        \end{enumerate}
    \end{latihan}
    \onslide<2->{%
        \begin{block}{Cek jawaban}
            \begin{enumerate}
                \item Konvergen ke $1/3$ karena $e^{-3R}\to0$.
                \item Divergen. Limit di sebelah kiri $1$ adalah $-\infty$
                    dan di sebelah kanan adalah $+\infty$.
                \item Pisahkan di $1$. Syarat $p<1$ di nol dan $p>1$
                    di tak hingga tidak dapat dipenuhi bersamaan.
            \end{enumerate}
        \end{block}
    }
\end{frame}

% KALKULUS 2 BAB 3.pdf, halaman PDF 25--37 (cetak 131--143).
\section{Limit Bentuk Tak Tentu}

\begin{frame}{\insertsection}
    \framesubtitle{Bentuk limit yang sama dapat memberi hasil berbeda}
    \begin{definisi}[Bentuk tak tentu $0/0$]
        Pada limit hasil bagi $f(x)/g(x)$, bentuk $0/0$ berarti
        $f(x)\to0$ dan $g(x)\to0$. Kedua informasi ini belum menentukan
        limit hasil baginya.
    \end{definisi}
    \onslide<2->{%
        \[
            \lim_{x\to0^+}\frac{x^2}{x}=0,\qquad
            \lim_{x\to0^+}\frac{3x}{x}=3,\qquad
            \lim_{x\to0^+}\frac{x}{x^2}=+\infty.
        \]
        Pada ketiga contoh, pembilang dan penyebut sama-sama mendekati nol.
        Laju masing-masing mendekati nol menghasilkan limit yang berbeda.
    }
\end{frame}

\begin{frame}{\insertsection}
    \framesubtitle{Aturan L'Hôpital untuk hasil bagi $0/0$}
    \begin{teorema}[Bentuk $0/0$]
        Misalkan $f,g\to0$ saat $x\to a$, termasuk limit satu sisi
        dan limit menuju $\pm\infty$. Anggap $f,g$ dapat diturunkan dan
        $g'\ne0$ pada selang yang didekati, kecuali mungkin di titik limit.
        Jika $\lim_{x\to a} f'(x)/g'(x)=L$ dengan $L$ berhingga atau $\pm\infty$, maka
        \[
            \lim_{x\to a}\frac{f(x)}{g(x)}
              =\lim_{x\to a}\frac{f'(x)}{g'(x)}=L.
        \]
    \end{teorema}
    \onslide<2->{Untuk $a$ berhingga, perpanjang $f(a)=g(a)=0$.
    Teorema Nilai Rata-rata Cauchy memberi
    $f(x)/g(x)=f'(\xi)/g'(\xi)$ dengan $\xi$ di antara $a$ dan $x$.
    Ketika $x\to a$, $\xi$ juga mendekati $a$.}
\end{frame}

\begin{frame}{\insertsection}
    \framesubtitle{Turunkan pembilang dan penyebut secara terpisah}
    \begin{columns}[t,onlytextwidth]
        \begin{column}{.47\textwidth}
            \begin{contoh}
                $\displaystyle\lim_{x\to3}\frac{x^2-9}{x-3}$
            \end{contoh}
            \onslide<2->{Kedua bagian menuju nol. Dengan L'Hôpital,
            \[
                \lim_{x\to3}\frac{2x}{1}=6.
            \]
            Hasil ini sama dengan limit $x+3$ setelah faktornya disederhanakan.}
        \end{column}
        \begin{column}{.49\textwidth}
            \begin{contoh}
                $\displaystyle\lim_{x\to0}\frac{\sin(3x)}{\sin(5x)}$
            \end{contoh}
            \onslide<3->{Bentuknya $0/0$. Ambil turunan pada kedua bagian.
            \[
                \lim_{x\to0}\frac{3\cos(3x)}{5\cos(5x)}=\frac35.
            \]
            Faktor $3$ dan $5$ muncul dari aturan rantai.}
        \end{column}
    \end{columns}
    \onslide<3->{L'Hôpital memakai hasil bagi $f'/g'$. Aturan turunan hasil
    bagi menghitung suatu besaran lain, yaitu $(f/g)'$.}
\end{frame}

\begin{frame}{\insertsection}
    \framesubtitle{Periksa bentuknya lagi sebelum mengulang}
    \begin{contoh}
        Hitung $\displaystyle\lim_{x\to0}\frac{1-\cos(2x)}{x^2}$.
    \end{contoh}
    \onslide<2->{Bentuk awalnya $0/0$, jadi ambil turunan pembilang dan penyebut.
    \[
        \lim_{x\to0}\frac{1-\cos(2x)}{x^2}
        =\lim_{x\to0}\frac{2\sin(2x)}{2x}.
    \]}
    \onslide<3->{Hasil antara masih $0/0$. Terapkan L'Hôpital sekali lagi.
    \[
        \lim_{x\to0}\frac{4\cos(2x)}{2}=2.
    \]
    Pemeriksaan tiap putaran menentukan kapan diferensiasi dapat dihentikan.}
\end{frame}

\begin{frame}{\insertsection}
    \framesubtitle{Bentuk $0/0$ juga muncul pada limit satu sisi dan tak hingga}
    \begin{columns}[t,onlytextwidth]
        \begin{column}{.48\textwidth}
            \begin{contoh}
                $\displaystyle\lim_{x\to0^\pm}\frac{\tan(2x)}{x^2}$
            \end{contoh}
            \onslide<2->{L'Hôpital pada masing-masing sisi memberi limit
            $\sec^2(2x)/x$. Dari kanan hasilnya $+\infty$, dari kiri
            $-\infty$. Limit dua sisi tidak ada.}
        \end{column}
        \begin{column}{.48\textwidth}
            \begin{contoh}
                $\displaystyle\lim_{x\to\infty}\frac{x^{-3/2}}{\sin(1/x)}$
            \end{contoh}
            \onslide<3->{Saat $x$ membesar, keduanya menuju nol.
            \[
                \begin{aligned}
                    \lim_{x\to\infty}\frac{-(3/2)x^{-5/2}}{-x^{-2}\cos(1/x)}
                        &=\lim_{x\to\infty}\frac{3}{2\sqrt{x}\cos(1/x)}\\
                        &=0.
                \end{aligned}
            \]}
        \end{column}
    \end{columns}
    \onslide<3->{Periksa perilaku fungsi pada arah limit yang diminta.}
\end{frame}

\begin{frame}{\insertsection}
    \framesubtitle{Kapan langsung memakai tanda atau apitan?}
    \begin{columns}[t,onlytextwidth]
        \begin{column}{.47\textwidth}
            \begin{contoh}
                $\displaystyle\lim_{x\to0}\frac{e^{2x}}{x^2}$
            \end{contoh}
            \onslide<2->{Pembilang menuju $1$ dan penyebut menuju $0$ dari
            nilai positif. Karena $e^{2x}\ge 1/2$ untuk $x$ cukup dekat ke $0$,
            hasil baginya menuju $+\infty$.}
        \end{column}
        \begin{column}{.49\textwidth}
            \begin{contoh}
                $\displaystyle\lim_{x\to\infty}\frac{x+\sin x}{x}$
            \end{contoh}
            \onslide<3->{Hasil bagi turunan $1+\cos x$ berosilasi.
            Gunakan $\dfrac{x+\sin x}{x}=1+\dfrac{\sin x}{x}$.
            Karena $|\sin x|/x\le1/x$ untuk $x>0$, limitnya $1$.}
        \end{column}
    \end{columns}
    \onslide<3->{Jika syarat L'Hôpital gagal, bentuk atau batas fungsi
    semula masih dapat menentukan limitnya.}
\end{frame}

\begin{frame}{\insertsection}
    \framesubtitle{Aturan L'Hôpital untuk hasil bagi tak hingga}
    \begin{teorema}[Bentuk $\infty/\infty$]
        Jika $f(x)$ dan $g(x)$ menuju $+\infty$ atau $-\infty$,
        gunakan syarat keterdiferensialan dan $g'\ne0$ seperti pada kasus $0/0$.
        Apabila $\lim f'/g'=L$ berhingga atau $\pm\infty$, maka $\lim f/g=L$.
    \end{teorema}
    \begin{contoh}
        Untuk $x\to\infty$, baik $\ln x$ maupun $\sqrt{x}$ menuju $+\infty$.
        \onslide<2->{%
            \[
                \lim_{x\to\infty}\frac{\ln x}{\sqrt{x}}
                =\lim_{x\to\infty}\frac{1/x}{1/(2\sqrt{x})}
                =\lim_{x\to\infty}\frac{2}{\sqrt{x}}=0.
            \]
            Pangkat $\sqrt{x}$ tumbuh lebih cepat daripada $\ln x$.}
    \end{contoh}
\end{frame}

\begin{frame}{\insertsection}
    \framesubtitle{Bentuk $0\cdot\infty$ dapat diubah menjadi hasil bagi}
    Jika $f\to0$ dan $|g|\to\infty$, tulis $fg=f/(1/g)$ atau $g/(1/f)$
    pada selang yang sesuai. Pilih bentuk dengan turunan yang lebih mudah.
    \begin{columns}[t,onlytextwidth]
        \begin{column}{.48\textwidth}
            \begin{contoh}
                $\displaystyle\lim_{x\to0^+}x^2\ln x$
            \end{contoh}
            \onslide<2->{Tulis $\ln x/x^{-2}$. Bentuk tak hingga per tak hingga memberi
            \[
                \lim_{x\to0^+}\frac{1/x}{-2x^{-3}}
                  =\lim_{x\to0^+}-\frac{x^2}{2}=0.
            \]}
        \end{column}
        \begin{column}{.48\textwidth}
            \begin{contoh}
                $\displaystyle\lim_{x\to\infty}xe^{-2x}$
            \end{contoh}
            \onslide<3->{Tulis $x/e^{2x}$. Kini kedua bagian menuju tak hingga.
            \[
                \lim_{x\to\infty}\frac{1}{2e^{2x}}=0.
            \]}
        \end{column}
    \end{columns}
\end{frame}

\begin{frame}{\insertsection}
    \framesubtitle{Bentuk $\infty-\infty$ dapat disatukan penyebutnya}
    \begin{contoh}
        Tentukan $\displaystyle\lim_{x\to0^+}
        \left(\frac1x-\frac1{\ln(1+x)}\right)$.
    \end{contoh}
    \onslide<2->{Kedua suku menuju $+\infty$. Samakan penyebut agar terbentuk $0/0$.
    \[
        \frac1x-\frac1{\ln(1+x)}
        =\frac{\ln(1+x)-x}{x\ln(1+x)}.
    \]}
    \onslide<3->{Gunakan L'Hôpital. Hasil antara masih berbentuk $0/0$, sehingga
    dapat diulangi.
    \[
        \lim_{x\to0^+}\frac{-x/(1+x)}{\ln(1+x)+x/(1+x)}
        =\lim_{x\to0^+}\frac{-1/(1+x)^2}{1/(1+x)+1/(1+x)^2}
        =-\frac12.
    \]}
\end{frame}

\begin{frame}{\insertsection}
    \framesubtitle{Sekawan menyingkap selisih dua suku besar}
    \begin{columns}[t,onlytextwidth]
        \begin{column}{.53\textwidth}
            \begin{contoh}
                $\displaystyle\lim_{x\to\infty}(\sqrt{x^2+6x}-x)$
            \end{contoh}
            \onslide<2->{Kalikan dengan bentuk sekawan.
            \[
                \frac{6x}{\sqrt{x^2+6x}+x}
                =\frac{6}{\sqrt{1+6/x}+1}\longrightarrow3.
            \]}
        \end{column}
        \begin{column}{.43\textwidth}
            \begin{block}{Perhatikan tanda tiap suku}
                Saat $x\to0^+$, $1/x\to+\infty$ dan $\ln x\to-\infty$.
                Jadi $1/x-\ln x$ adalah penjumlahan dua kontribusi positif
                yang membesar. Limitnya $+\infty$.
            \end{block}
        \end{column}
    \end{columns}
\end{frame}

\begin{frame}{\insertsection}
    \framesubtitle{Bentuk pangkat ditangani dengan logaritma}
    \begin{definisi}[Bentuk pangkat tak tentu]
        Untuk $y=f(x)^{g(x)}$ dengan $f(x)>0$ di dekat titik limit,
        bentuk $0^0$, $\infty^0$, dan $1^\infty$ menyatakan pasangan
        limit basis dan pangkat yang belum menentukan limit $y$.
    \end{definisi}
    \[
        \ln y=g(x)\ln f(x),\qquad y=\exp\bigl(g(x)\ln f(x)\bigr).
    \]
    \onslide<2->{Cari limit $\ln y$ melalui bentuk perkalian atau hasil bagi.
    Jika limit $\ln y=L$ berhingga, maka $y\to e^L$.
    Jika $\ln y\to+\infty$, maka $y\to+\infty$.
    Jika $\ln y\to-\infty$, maka $y\to0$.}
\end{frame}

\begin{frame}{\insertsection}
    \framesubtitle{Dua pangkat yang sama-sama menuju nol}
    \begin{columns}[t,onlytextwidth]
        \begin{column}{.47\textwidth}
            \begin{contoh}
                $\displaystyle\lim_{x\to0^+}(2x)^x$
            \end{contoh}
            \onslide<2->{Bentuknya $0^0$. Logaritmanya
            $x\ln(2x)=x\ln2+x\ln x$. Karena $x\ln x\to0$,
            maka $\ln y\to0$ dan $y\to1$.}
        \end{column}
        \begin{column}{.49\textwidth}
            \begin{contoh}
                $\displaystyle\lim_{x\to\infty}(\sqrt{x})^{\,2/(1+\ln x)}$
            \end{contoh}
            \onslide<3->{Bentuknya $\infty^0$. Logaritmanya
            \[
                \frac{2\ln(\sqrt{x})}{1+\ln x}
                  =\frac{\ln x}{1+\ln x}\longrightarrow1.
            \]
            Jadi $y\to e$.}
        \end{column}
    \end{columns}
\end{frame}

\begin{frame}{\insertsection}
    \framesubtitle{Bentuk $1^\infty$ dan limit yang menghasilkan $e$}
    \begin{contoh}
        Cari $\displaystyle\lim_{x\to0^+}(1+3x)^{2/x}$.
    \end{contoh}
    \onslide<2->{Ambil logaritma. Bentuknya sekarang $0/0$.
    \[
        \lim_{x\to0^+}\ln y
         =\lim_{x\to0^+}\frac{2\ln(1+3x)}{x}
         =\lim_{x\to0^+}\frac6{1+3x}=6.
    \]}
    \onslide<3->{Maka limit asalnya $e^6$.
    Dengan cara yang sama, $\ln(1+x)/x\to1$ dari kedua sisi, sehingga
    $(1+x)^{1/x}\to e$ saat $x\to0$.}
\end{frame}

\begin{frame}{\insertsection}
    \framesubtitle{Membandingkan dua luas yang menyusut}
    \begin{columns}[c,onlytextwidth]
        \begin{column}{.56\textwidth}
            \begin{contoh}
                Pada lingkaran berjari-jari $2$, bandingkan luas $T$ dari
                segitiga $ABC$ dengan luas segmen $S$ antara tali busur $AB$
                dan busur $AB$ saat $\theta\to0^+$.
            \end{contoh}
            \onslide<2->{%
                $AC=2\sin\theta$ dan $CB=2(1-\cos\theta)$.
                Luas segmen diperoleh dari juring $OAB$ dikurangi segitiga $OAB$.
                \[
                    \frac{T}{S}
                    =\frac{2\sin\theta(1-\cos\theta)}{2(\theta-\sin\theta)}.
                \]}
            \onslide<3->{Ini bentuk $0/0$. Hasil bagi turunannya
            $\cos\theta+\sin^2\theta/(1-\cos\theta)=1+2\cos\theta$.
            Maka $\lim_{\theta\to0^+}T/S=3$.}
        \end{column}
        \begin{column}{.40\textwidth}
            \centering
            \begin{tikzpicture}
                \coordinate (O) at (0,0);
                \coordinate (B) at (3.45,0);
                \coordinate (A) at ({3.45*cos(53)},{3.45*sin(53)});
                \coordinate (C) at ({3.45*cos(53)},0);
                \fill[HIMAmuda!36] (A) -- (C) -- (B) -- cycle;
                \fill[magenta!24] (B) arc[start angle=0,end angle=53,radius=3.45] -- cycle;
                \draw[HIMAtua,very thick] (B) arc[start angle=0,end angle=53,radius=3.45];
                \draw[HIMAtua,thick] (O) -- (A) -- (B) -- cycle;
                \draw[HIMAtua,thick] (A) -- (C);
                \draw (C) ++(.17,0) -- ++(0,.17) -- ++(-.17,0);
                \draw[teal!80!black] (.69,0) arc[start angle=0,end angle=53,radius=.69];
                \node at (.87,.37) {$\theta$};
                \node[below left] at (O) {$O$};
                \node[above] at (A) {$A$};
                \node[below] at (B) {$B$};
                \node[below] at (C) {$C$};
                \node[teal!80!black] at (2.31,.53) {$T$};
                \node[magenta!90!black] at (2.97,1.52) {$S$};
            \end{tikzpicture}
        \end{column}
    \end{columns}
\end{frame}

\begin{frame}{\insertsection}
    \framesubtitle{Limit logaritmik menentukan nilai integral tak wajar}
    \begin{contoh}
        Hitung $\displaystyle\int_0^1 x\ln x\,dx$.
    \end{contoh}
    \onslide<2->{Pada $[s,1]$ dengan $s>0$, gunakan integrasi parsial
    dengan $u=\ln x$ dan $dv=x\,dx$.
    \[
        \int_s^1 x\ln x\,dx
           =\left[\frac{x^2}{2}\ln x-\frac{x^2}{4}\right]_s^1.
    \]}
    \onslide<3->{Dari $s^2\ln s\to0$ saat $s\to0^+$, suku di batas bawah
    menuju nol. Maka integralnya konvergen ke $-1/4$.
    Tanda negatif sesuai dengan $x\ln x<0$ pada $(0,1)$.}
\end{frame}

\begin{frame}{\insertsection}
    \framesubtitle{Fungsi Gamma diuji di dua ujung}
    \begin{definisi}[Fungsi Gamma]
        Untuk $s>0$, fungsi Gamma didefinisikan oleh
        \[
            \Gamma(s)=\int_0^\infty t^{s-1}e^{-t}\,dt.
        \]
    \end{definisi}
    \begin{columns}[t,onlytextwidth]
        \begin{column}{.48\textwidth}
            \onslide<2->{\textbf{Dekat nol.}
            Pada $(0,1]$, $e^{-1}\le e^{-t}\le1$.
            Maka konvergensi mengikuti $\int_0^1 t^{s-1}\,dt$,
            yang konvergen tepat saat $s>0$.}
        \end{column}
        \begin{column}{.48\textwidth}
            \onslide<3->{\textbf{Dekat tak hingga.}
            $t^{s+1}e^{-t}\to0$ karena eksponensial mengalahkan pangkat.
            Jadi $t^{s-1}e^{-t}\le t^{-2}$ untuk $t$ cukup besar, dan
            ekornya konvergen.}
        \end{column}
    \end{columns}
\end{frame}

\begin{frame}{\insertsection}
    \framesubtitle{Integrasi parsial menurunkan pangkat pada Gamma}
    \begin{teorema}[Relasi rekursif Gamma]
        Untuk $s>0$, $\Gamma(s+1)=s\Gamma(s)$ dan $\Gamma(1)=1$.
    \end{teorema}
    \onslide<2->{Integrasikan $t^se^{-t}$ secara parsial pada $[\varepsilon,R]$.
    \[
        \int_\varepsilon^R t^se^{-t}\,dt
          =[-t^se^{-t}]_\varepsilon^R+
            s\int_\varepsilon^R t^{s-1}e^{-t}\,dt.
    \]
    Suku batas menuju nol di kedua ujung. Untuk $s=1$ pada definisi
    Gamma, $\int_0^\infty e^{-t}\,dt=1$.}
    \onslide<3->{%
        \begin{contoh}
            $\Gamma(4)=3\Gamma(3)=3\cdot2\cdot1=6$.
            Substitusi $t=x^2$ pada $\Gamma(1/2)$ memberi
            $2\int_0^\infty e^{-x^2}\,dx=\sqrt{\pi}$.
        \end{contoh}
    }
\end{frame}

\begin{frame}{\insertsection}
    \framesubtitle{Latihan limit bentuk tak tentu}
    \begin{latihan}
        Tentukan bentuk awal dan nilai limit.
        \begin{enumerate}
            \item $\displaystyle\lim_{x\to0}\frac{\sin(4x)}{e^{2x}-1}$.
            \item $\displaystyle\lim_{x\to\infty}\frac{x^2}{e^x}$.
            \item $\displaystyle\lim_{x\to\infty}(1-2/x)^{3x}$.
        \end{enumerate}
    \end{latihan}
    \onslide<2->{%
        \begin{block}{Cek jawaban}
            \begin{enumerate}
                \item $0/0$. Hasil bagi turunan menuju $4/2=2$.
                \item $\infty/\infty$. Dua penerapan L'Hôpital memberi
                    limit $2/e^x\to0$.
                \item $1^\infty$. Logaritmanya $3\ln(1-2/x)/(1/x)\to-6$,
                    sehingga hasilnya $e^{-6}$.
            \end{enumerate}
        \end{block}
    }
\end{frame}

\end{document}
